\documentclass[11pt]{article}
\usepackage[english]{babel}
\usepackage[a4paper,margin=0.9in]{geometry}
\usepackage[T1]{fontenc}
\usepackage[utf8]{inputenc}
\usepackage{lmodern}
\usepackage{amsmath,amssymb,mathtools}
\usepackage{microtype}
\usepackage{fancyhdr}
\usepackage{array,enumitem}
\setlength{\parindent}{1.4em}
\setlength{\parskip}{0.15em}
\emergencystretch=6em
\setlength{\headheight}{14pt}
\pagestyle{fancy}
\fancyhf{}
\cfoot{\thepage}
\newcommand{\sect}[2]{\section*{\S\,#1. #2}}
\newcommand{\D}{\mathrm D}

\lhead{Weber, Lehrbuch der Algebra I}
\rhead{English translation: \S\,55--\S\,61}
\begin{document}
\begin{center}
{\large Heinrich Weber}\\[0.35em]
{\Large Textbook of Algebra. Volume I.}\\[0.35em]
{\large Fifth Section. Linear Transformation.}\\[0.25em]
{\large \S\,55--\S\,61}
\end{center}

\sect{55}{Introduction of the linear transformation}

When, in any functions depending on the $m$ variables $x_1,x_2,\ldots,x_m$, we substitute for these variables linear expressions depending on $m$ new variables $x'_1,x'_2,\ldots,x'_m$,
\begin{equation}
\begin{aligned}
 x_1&=\alpha_1^{(1)}x'_1+\alpha_1^{(2)}x'_2+\cdots+\alpha_1^{(m)}x'_m,\\
 x_2&=\alpha_2^{(1)}x'_1+\alpha_2^{(2)}x'_2+\cdots+\alpha_2^{(m)}x'_m,\\
 &\hspace{3.0em}\cdots\\
 x_m&=\alpha_m^{(1)}x'_1+\alpha_m^{(2)}x'_2+\cdots+\alpha_m^{(m)}x'_m,
\end{aligned}
\tag{1}
\end{equation}
we call this a linear substitution, and the result a linear transformation.

The quantity
\begin{equation}
 r=\begin{vmatrix}
 \alpha_1^{(1)}&\alpha_1^{(2)}&\cdots&\alpha_1^{(m)}\\
 \alpha_2^{(1)}&\alpha_2^{(2)}&\cdots&\alpha_2^{(m)}\\
 \vdots&\vdots&&\vdots\\
 \alpha_m^{(1)}&\alpha_m^{(2)}&\cdots&\alpha_m^{(m)}
 \end{vmatrix}
\tag{2}
\end{equation}
is called the substitution determinant. It must always be different from zero, since otherwise (1) would express a dependence among the variables $x_1,x_2,\ldots,x_m$.

Under a linear substitution, every homogeneous function of the variables $x$ passes into a homogeneous function of the same degree in the variables $x'$.

Thus, for example, every linear function
\[
 y=\sum_i a_i x_i
\]
passes into a similar function
\[
 y=\sum_i a'_i x'_i,
\]
where
\begin{equation}
 a'_k=\sum_i a_i\alpha_i^{(k)}.
\tag{3}
\end{equation}
If we consider $m$ such linear functions
\[
 y_\nu=\sum_i a_{\nu,i}x_i\qquad \nu=1,2,\ldots,m,
\]
we obtain $m$ transformed functions
\[
 y_\nu=\sum_i a'_{\nu,i}x'_i\qquad \nu=1,2,\ldots,m.
\]
The determinants of the systems $y_\nu,y'_\nu$,
\[
 \Delta=\sum \pm a_{1,1}a_{2,2}\cdots a_{m,m},
 \qquad
 \Delta'=\sum \pm a'_{1,1}a'_{2,2}\cdots a'_{m,m},
\]
stand, by the multiplication theorem for determinants, in the relation
\begin{equation}
 \Delta'=r\Delta,
\tag{4}
\end{equation}
so that one cannot vanish without the other.

\sect{56}{Quadratic forms}

Of particular importance is the transformation of homogeneous functions of the second degree, or quadratic forms. We denote them, as already in the first section, by
\begin{equation}
 \varphi(x_1,x_2,\ldots,x_m)=\sum_{i,k}a_{i,k}x_i x_k,
 \qquad a_{i,k}=a_{k,i}.
\tag{1}
\end{equation}
By the substitution \S\,55, (1), $\varphi$ passes into
\begin{equation}
 \Phi(x'_1,x'_2,\ldots,x'_m)=\sum_{i,k}a'_{i,k}x'_i x'_k.
\tag{2}
\end{equation}
Put $x+\xi$ for $x$, and similarly $x'+\xi'$ for $x'$, so that the $\xi'$ stand in the same dependence on the $\xi$ as the $x'$ do on the $x$. If, for abbreviation, we set
\[
 \frac12\varphi'(x_i)=u_i,
 \qquad
 \frac12\Phi'(x'_i)=u'_i,
\]
then, by comparison of the terms of the first dimension (\S\,16), we obtain
\begin{equation}
 \sum \xi_i u_i=\sum \xi'_i u'_i,
\tag{3}
\end{equation}
and hence, as in \S\,55, (3),
\begin{equation}
 u'_k=\sum_i \alpha_i^{(k)}u_i.
\tag{4}
\end{equation}
Now
\begin{equation}
 u_k=\sum_i a_{k,i}x_i,
 \qquad
 u'_k=\sum_i a'_{k,i}x'_i,
\tag{5}
\end{equation}
and consequently (4) gives us the transformation of $m$ linear functions. The coefficients of the given system are $a'_{k,h}$, and those of the transformed system, obtained by substituting the first system (5) into (4), are
\begin{equation}
 b_{k,h}=\sum_i \alpha_i^{(k)}a_{i,h}.
\tag{6}
\end{equation}
The determinant of the first system
\[
 H'=\sum \pm a'_{1,1}a'_{2,2}\cdots a'_{m,m}
\]
is, by \S\,55, (4),
\[
 =r\sum\pm b_{1,1}b_{2,2}\cdots b_{m,m},
\]
and, by (6), the determinant of the $b_{k,h}$ is
\[
 =r\sum\pm a_{1,1}a_{2,2}\cdots a_{m,m}.
\]
If we therefore put
\begin{equation}
 H=\sum \pm a_{1,1}a_{2,2}\cdots a_{m,m},
\tag{7}
\end{equation}
we have the important theorem
\begin{equation}
 H'=r^2H.
\tag{8}
\end{equation}
The determinant $H$ is called the determinant of the form $\varphi(x)$.

It therefore changes, under the linear transformation of the function $\varphi$, only by the factor $r^2$, and it cannot vanish for the transformed form if it does not vanish for the original one.

The vanishing of the determinant $H$ means that the function $\varphi(x)$ may be regarded as a function of fewer than $m$ variables that are linear functions of the $x$.

For, first, if by a linear transformation $\varphi(x)$ passes into the function $\Phi(x')$, which is independent of one of the variables, say $x'_1$, then $H'$ vanishes, since all the elements in one of the rows and columns vanish; consequently, by (8), $H=0$ as well.

Second, however, $H$ is the determinant of the system of linear equations
\begin{equation}
 \varphi'(x_i)=0\qquad i=1,2,\ldots,m,
\tag{9}
\end{equation}
and, if this determinant vanishes, then (\S\,24, III) there is a system of quantities
\begin{equation}
 x_1^{(0)},x_2^{(0)},\ldots,x_m^{(0)},
\tag{10}
\end{equation}
whose elements do not all vanish and which, when substituted for the $x$, satisfy equations (9).

Now if $x_i,x'_i$ are any two systems of quantities, and $\lambda$ is likewise an arbitrary quantity, the following identities hold:
\begin{equation}
\begin{aligned}
 \varphi(x+\lambda x')&=\varphi(x)+\lambda\sum x_i\varphi'(x'_i)+\lambda^2\varphi(x'),\\
 \varphi(x')&=\frac12\sum x'_i\varphi'(x'_i),
\end{aligned}
\tag{11}
\end{equation}
[\S\,16, (12) to (15)]. Thus, when for $x'_i$ one substitutes the values $x_i^{(0)}$, one obtains
\begin{equation}
 \varphi(x_i+\lambda x_i^{(0)})=\varphi(x),
\tag{12}
\end{equation}
quite independently of $\lambda$. If we now assume that $x_1^{(0)}$ is different from zero and set
\[
 \lambda=-\frac{x_1}{x_1^{(0)}},
\]
and further set
\begin{equation}
\begin{aligned}
 y_2&=x_2x_1^{(0)}-x_1x_2^{(0)},\\
 y_3&=x_3x_1^{(0)}-x_1x_3^{(0)},\\
 &\hspace{3em}\cdots\\
 y_m&=x_mx_1^{(0)}-x_1x_m^{(0)},
\end{aligned}
\tag{13}
\end{equation}
then it follows that
\begin{equation}
 x_1^{(0)2}\varphi(x)=\varphi(0,y_2,y_3,\ldots,y_m),
\tag{14}
\end{equation}
so that $\varphi(x)$ is expressed through the $m-1$ variables $y_2,y_3,\ldots,y_m$.

\sect{57}{Transformation of quadratic forms into a sum of squares}

We at once make an important application of the theorem just proved to the discriminant of the function (11), regarded as a function of the second degree in $\lambda$, without now assuming that the determinant of the function $\varphi$ vanishes. This discriminant is
\begin{equation}
 \psi(x)=\left[\sum x_i\varphi'(x'_i)\right]^2-4\varphi(x)\varphi(x').
\tag{1}
\end{equation}
If, in this expression, we put $x_i+\lambda x'_i$ in place of $x_i$, then, after substituting for $\varphi(x+\lambda x')$ the expression \S\,56, (11) and setting $\sum x'_i\varphi'(x'_i)=2\varphi(x')$, we get
\[
 \psi(x+\lambda x')=\psi(x),
\]
which shows that $\psi(x)$ depends on only $m-1$ linear combinations of the $m$ variables. If for $x'_1,x'_2,\ldots,x'_m$ we put any special values for which $\varphi(x')$ does not vanish, and assume that $x'_1$ is not zero, then
\[
 x_1'^2\psi(x_1,x_2,\ldots,x_m)=
 \psi(0,x_2x'_1-x_1x'_2,\ldots,x_mx'_1-x_1x'_m),
\]
hence it is a function of the $m-1$ variables
\[
 y_2=x_2x'_1-x_1x'_2,
 \ldots,
 y_m=x_mx'_1-x_1x'_m.
\]
This may also be confirmed by the fact that
\[
 \psi'(x_1),\psi'(x_2),\ldots,\psi'(x_m)
\]
all vanish simultaneously for $x_1,x_2,\ldots,x_m=x'_1,x'_2,\ldots,x'_m$, so that the determinant of the function $\psi$ must vanish.

Formula (1) now contains a very important result, most clearly seen when we give this formula the following form. We set
\begin{equation}
 \psi(x)=-4\varphi(x')\chi(y_2,\ldots,y_m),
 \qquad
 \sum x_i\varphi'(x'_i)=2\sqrt{\varphi(x')}\,Y_1,
\tag{2}
\end{equation}
and obtain
\begin{equation}
 \varphi(x)=Y_1^2+\chi(y_2,y_3,\ldots,y_m),
\tag{3}
\end{equation}
so that $\varphi(x)$ is represented as the sum of the square of a linear function $Y_1$ and a function $\chi$ of $m-1$ variables.

From this follows the important theorem: A quadratic form in $m$ variables can always be represented as a sum of squares of $m$ or fewer linear functions.

This theorem is obviously true for a function of one variable, since such a function is itself a square. From (3), however, the truth of the theorem for a function of $m$ variables follows if its truth is assumed for a function of $m-1$ variables.

As (2) shows, the representation is possible in infinitely many different ways, because the entirely arbitrary quantities $x'$ still occur in this formula.

A representation by fewer than $m$ squares is possible only when the determinant of the function $\varphi$ vanishes.

Since in the expression (2) for $Y_1$ a square root occurs, $Y_1$ may be real or imaginary even when the coefficients of $\varphi$ and the $x$ and $x'$ are assumed real. But if, in the case where $Y_1$ is imaginary, we put $iY_1$ in place of $Y_1$, then we can also state our theorem as follows:

A real quadratic form in $m$ variables can be represented as a sum of at most $m$ positive or negative squares of real linear functions of the $x$.\footnote{For the general theory of the linear transformation of forms of the second degree, besides the investigations of Jacobi and Hesse (Hesse, ``Vorlesungen über analytische Geometrie des Raumes'', third edition, edited by Gundelfinger, Leipzig 1876), especially the works of Weierstrass and Kronecker (Monthly Reports of the Berlin Academy, 1868, 1874) are important. A detailed historical account of the development of the question is given in the ``Bericht über den gegenwärtigen Stand der Invariantentheorie'' by Franz Meyer (first annual report of the Deutsche Mathematiker-Vereinigung, Berlin 1892).}

\sect{58}{Law of inertia of quadratic forms}

For the decomposition of a quadratic form into squares, the following theorem now holds. Sylvester gave it the name of the law of inertia of quadratic forms (Philos. Mag. 1852).

In whatever way one decomposes a real quadratic form $\varphi(x)$ into a sum of positive and negative squares of linear functions, the number of positive and negative squares, and therefore also their total number, is always the same, provided that there is no linear dependence among these linear functions.

The proof of this important theorem is very simple.

Let
\begin{equation}
\begin{aligned}
 \varphi(x)&=Y_1^2+Y_2^2+\cdots+Y_\nu^2-Y_1'^2-Y_2'^2-\cdots-Y_{\nu'}'^2\\
 &=Z_1^2+Z_2^2+\cdots+Z_\mu^2-Z_1'^2-Z_2'^2-\cdots-Z_{\mu'}'^2
\end{aligned}
\tag{1}
\end{equation}
be two decompositions of $\varphi(x)$ into squares. Thus $Y_1,Y_2,\ldots,Y_\nu,Y'_1,Y'_2,\ldots,Y'_{\nu'}$ are homogeneous linear functions of $x$ among which no linear relation with constant coefficients exists, so that $\nu+\nu'\leq m$. The same assumptions are made about the functions $Z_1,Z_2,\ldots,Z_\mu,Z'_1,Z'_2,\ldots,Z'_{\mu'}$, so that $\mu+\mu'\leq m$ as well.

Suppose that
\[
 \nu<\mu.
\]
Then we can subject the variables $x$ to the linear equations
\begin{equation}
\begin{aligned}
 Y_1&=0,\quad Y_2=0,\ldots,Y_\nu=0,\\
 Z'_1&=0,\quad Z'_2=0,\ldots,Z'_{\mu'}=0,
\end{aligned}
\tag{2}
\end{equation}
whose number $\nu+\mu'$ is less than $m$.

If all the functions $Z_1,Z_2,\ldots,Z_\mu$ were linearly dependent on the functions $Y_1,Y_2,\ldots,Y_\nu,Z'_1,Z'_2,\ldots,Z'_{\mu'}$, then from these $\mu$ equations one could, by eliminating the $\nu$ variables $Y_1,Y_2,\ldots,Y_\nu$, derive a linear equation among $Z_1,Z_2,\ldots,Z_\mu,Z'_1,\ldots,Z'_{\mu'}$, which by hypothesis is not to exist. Thus the vanishing of all $Z_1,Z_2,\ldots,Z_\mu$ is not a necessary consequence of equations (2), and we may add to equations (2) a non-homogeneous equation, for instance
\[
 Z_1=1.
\]
Then, for the $m$ unknowns $x$, we have a system of $m$ or fewer equations that are independent of one another and hence do not contradict one another.

But then the second representation (1) gives a positive value for $\varphi(x)$, while the first gives a negative or vanishing value, a contradiction. Hence
\[
 \nu\geq \mu,
\]
and, since one similarly concludes that
\[
 \mu\geq \nu,
\]
only $\nu=\mu$ remains. In the same way one may prove that $\nu'=\mu'$, and this proves the theorem completely.

\sect{59}{Transformation of forms of the $n$th degree}

When a form $F(x)$ of the $n$th degree in $m$ variables is transformed by a linear substitution [\S\,55, (1)] into $\Phi(x')$, the linear function of the variables $\xi_1,\xi_2,\ldots,\xi_m$
\begin{equation}
 \sum_i \xi_i F'(x_i)
\tag{1}
\end{equation}
is simultaneously transformed, by the same substitution applied to the $\xi_i$, into
\begin{equation}
 \sum_i \xi'_i\Phi'(x'_i),
\tag{2}
\end{equation}
because both expressions represent the coefficient of the first power of $t$ in the expansion of
\[
 F(x_i+t\xi_i)=\Phi(x'_i+t\xi'_i)
\]
in powers of $t$ (\S\,16). Likewise, the quadratic function
\begin{equation}
 \sum_{i,k}\xi_i\xi_k F''(x_i x_k)
\tag{3}
\end{equation}
is transformed into
\begin{equation}
 \sum_{i,k}\xi'_i\xi'_k\Phi''(x'_i x'_k).
\tag{4}
\end{equation}

If we apply the first result to a system of $m$ functions $F_1(x),F_2(x),\ldots,F_m(x)$ of arbitrary degrees and denote by $\Delta$ the determinant
\[
\Delta=\begin{vmatrix}
 F'_1(x_1)&F'_1(x_2)&\cdots&F'_1(x_m)\\
 F'_2(x_1)&F'_2(x_2)&\cdots&F'_2(x_m)\\
 \vdots&\vdots&&\vdots\\
 F'_m(x_1)&F'_m(x_2)&\cdots&F'_m(x_m)
\end{vmatrix},
\]
while $\Delta'$ denotes the corresponding determinant formed from the functions $\Phi_1,\Phi_2,\ldots,\Phi_m$, then by \S\,55, (4), again denoting the substitution determinant by $r$, we obtain
\begin{equation}
 \Delta'=r\Delta.
\tag{5}
\end{equation}
The function $\Delta$ is called the functional determinant of the $m$ functions $F_1,F_2,\ldots,F_m$.

Likewise, from the transformation (3), (4), if $H$ denotes the determinant
\[
H=\begin{vmatrix}
 F''(x_1,x_1)&F''(x_1,x_2)&\cdots&F''(x_1,x_m)\\
 F''(x_2,x_1)&F''(x_2,x_2)&\cdots&F''(x_2,x_m)\\
 \vdots&\vdots&&\vdots\\
 F''(x_m,x_1)&F''(x_m,x_2)&\cdots&F''(x_m,x_m)
\end{vmatrix}
\]
and $H'$ the corresponding determinant for the transformed function $\Phi$, we obtain by \S\,56, (8)
\begin{equation}
 H'=r^2H.
\tag{6}
\end{equation}
$H$ is called the Hessian determinant\footnote{Named after Hesse, who frequently used this determinant in geometrical investigations, especially in the theory of curves of the third order and in elimination problems, and recognized its significance.} of the function $F$. In both formulas (5) and (6), $r$ denotes the substitution determinant.

\sect{60}{Invariants and covariants}

In the linear transformation of homogeneous functions, certain functions of the coefficients always occur which, when the coefficients of the given form are replaced by the coefficients of the transformed form, change in no other way than by acquiring a factor that is a power of the substitution determinant. Such functions also occur under the simultaneous transformation of a system of several forms; they depend on the coefficients of all forms of the system. Examples are, for a system of $n$ linear forms, the determinant of the system, and, for a quadratic form, the determinant of that form.

These functions are called invariants of the form, or, when they depend on several forms, simultaneous invariants of the system.

Thus, if $I(a)$ is a function of the coefficients $a$ of a form $F(x)$, and if $a'$ are the coefficients of the transformed form $\Phi(x')$, then $I(a)$ is called an invariant if the identical relation
\begin{equation}
 I(a')=r^\lambda I(a)
\tag{1}
\end{equation}
holds. We shall call the exponent $\lambda$ the weight of the invariant. If this exponent is zero, then $I$ is called an absolute invariant.

One usually considers only integral rational invariants, that is, such functions $I$ as depend rationally on the coefficients $a$ and are moreover integral homogeneous functions of them. Simultaneous invariants are to be homogeneous with respect to the coefficients of each of the forms on which they depend.

Besides invariants, there are also functions that, in addition to the coefficients, contain the variables themselves and otherwise have the same properties as the invariants (1), namely the ``invariant property.'' Thus, if $C(x,a)$ is such a function, then
\begin{equation}
 C(x',a')=r^\lambda C(x,a).
\tag{2}
\end{equation}
Such functions are called covariants of the form $F(x)$. The functional determinant is a simultaneous covariant of a system of $m$ functions, and the Hessian determinant of a function of degree higher than the second is a covariant of a single form. Among covariants, too, one chiefly considers integral rational and homogeneous functions of $x$.

If $m$ is the number of variables $x$, $n$ the degree of the function $F(x)$, and $\nu$ and $\mu$ the degrees of $C(x,a)$ with respect to $x$ and to $a$, then the numbers $\lambda,\mu,\nu,m,n$ satisfy the relation
\begin{equation}
 n\mu=m\lambda+\nu.
\tag{3}
\end{equation}
This is proved by comparing the degree of the right and left sides of (2) with respect to the substitution coefficients. Namely, the $a'$ are of degree $n$ in the substitution coefficients, $r$ is of degree $m$, and $x$ is of degree one; this gives (3). The corresponding formula for invariants is obtained by putting $\nu=0$, just as invariants may in general be regarded as a special case of covariants.

We shall briefly discuss a general rule for the formation of invariants and covariants, which we have already used in the special cases of \S\,59.

If, according to \S\,16, (4), we arrange the function
\[
 F(x_1+t\xi_1,x_2+t\xi_2,
 \ldots,x_m+t\xi_m)
\]
by powers of $t$, then for the coefficient of $t^\nu$ we obtain
\begin{equation}
 F_\nu(x,\xi)=\frac1{\Pi(\nu)}
 \sum_{\alpha_1+\alpha_2+\cdots+\alpha_m=\nu}
 \frac{\Pi(\nu)}{\Pi(\alpha_1)\cdots\Pi(\alpha_m)}
 \xi_1^{\alpha_1}\xi_2^{\alpha_2}\cdots\xi_m^{\alpha_m}
 \D_{\alpha_1,\alpha_2,\ldots,\alpha_m}(F).
\tag{4}
\end{equation}
If we apply the linear transformation (1), \S\,55, simultaneously to the variables $\xi$ and $x$, then $F_\nu(x,\xi)$ passes into $\Phi_\nu(x',\xi')$, which is obtained from $F_\nu$ by replacing at the same time all letters $x,\xi,a$ by $x',\xi',a'$ and $F$ by $\Phi$. Hence:

I. If $F_\nu(x,\xi)$ is regarded as a form of the $\nu$th degree in $\xi$ and one forms an invariant of this form, whose coefficients therefore still depend on the $x$, one obtains a covariant of the form $F$.

Since the variables $x$ and $\xi$ undergo the same transformation, it follows similarly that:

II. If one forms a covariant of the function $F_\nu(x,\xi)$, regarded as a function of $\xi$, and then replaces the $\xi$ in it by the $x$, one obtains a covariant of the original function $F(x)$.

The same theorems also hold for simultaneous invariants and covariants.

For the functions $F_\nu(x,\xi)$ one has the theorem
\begin{equation}
 F_\nu(x,\xi)=F_{n-\nu}(\xi,x),
\tag{5}
\end{equation}
which follows by comparing corresponding powers of $t$ on both sides of the identity
\[
 F(x_1+t\xi_1,\ldots,x_m+t\xi_m)=t^nF(t^{-1}x_1+\xi_1,\ldots,t^{-1}x_m+\xi_m).
\]

The functions
\[
 P_\nu(x,\xi)=\frac{\Pi(\nu)\Pi(n-\nu)}{\Pi(n)}F_\nu(x,\xi)
\]
are called the polars of the function $F(x)$; more precisely, $P_\nu(x,\xi)$ is called the $\nu$th polar. They too satisfy the theorem
\[
 P_\nu(x,\xi)=P_{n-\nu}(\xi,x).
\]
The constant factor $\Pi(\nu)\Pi(n-\nu):\Pi(n)$ is added so that, when the form $F(x)$ is written with polynomial coefficients, the functions $P_\nu$ have coefficients as simple as possible.

For the binary form $f(x,y)$, to which in what follows we shall mainly apply these definitions, the following expressions for the polars result:
\[
 P_1(x,\xi)=\frac1n\{\xi f'(x)+\eta f'(y)\},
\]
\[
 P_2(x,\xi)=\frac1{n(n-1)}\{\xi^2f''(x,x)+2\xi\eta f''(x,y)+\eta^2f''(y,y)\},
\]
and in general
\[
\begin{aligned}
&n(n-1)\cdots(n-\nu+1)P_\nu(x,\xi)\\
&\quad=\xi^\nu\frac{\partial^\nu f}{\partial x^\nu}
+\nu\xi^{\nu-1}\eta\frac{\partial^\nu f}{\partial x^{\nu-1}\partial y}
+\cdots+
\eta^\nu\frac{\partial^\nu f}{\partial y^\nu}.
\end{aligned}
\]

\sect{61}{Linear transformation of binary forms}

An integral rational function of degree $n$ in $x$,
\begin{equation}
 f(x)=a_0x^n+a_1x^{n-1}+a_2x^{n-2}+\cdots+a_{n-1}x+a_n,
\tag{1}
\end{equation}
is converted into a binary form when $x$ is replaced by $x:y$ and the result is multiplied by $y^n$:
\begin{equation}
 f(x,y)=a_0x^n+a_1x^{n-1}y+a_2x^{n-2}y^2+\cdots+a_ny^n.
\tag{2}
\end{equation}
If the roots of $f(x)$ are known, then $f(x)$ can be decomposed into $n$ linear factors:
\begin{equation}
 f(x)=a_0(x-\alpha_1)(x-\alpha_2)\cdots(x-\alpha_n),
\tag{3}
\end{equation}
and from this it follows that
\begin{equation}
 f(x,y)=a_0(x-\alpha_1y)(x-\alpha_2y)\cdots(x-\alpha_ny).
\tag{4}
\end{equation}
But if $\alpha_1,\alpha_2,\ldots,\alpha_n$ are themselves replaced by ratios $\alpha_1:\beta_1,\alpha_2:\beta_2,\ldots,\alpha_n:\beta_n$, one obtains
\begin{equation}
 f(x,y)=A(x\beta_1-y\alpha_1)(x\beta_2-y\alpha_2)\cdots(x\beta_n-y\alpha_n),
\tag{5}
\end{equation}
provided that
\[
 a_0=A\beta_1\beta_2\cdots\beta_n.
\]

For the factors of $f(x,y)$ that will now occur more frequently, we introduce the abbreviated notation
\begin{equation}
 x\beta_1-y\alpha_1=(x\beta_1).
\tag{6}
\end{equation}
Similarly, we shall put
\begin{equation}
 x_1y_2-x_2y_1=(x_1y_2),
\tag{7}
\end{equation}
where $x_1,y_1$ and $x_2,y_2$ are two arbitrary pairs of variables. If we use the inhomogeneous representation, we shall write still more briefly
\begin{equation}
\begin{gathered}
 x-\alpha_1=(0,1),\quad x-\alpha_2=(0,2),\ldots,x-\alpha_n=(0,n),\\
 \alpha_1-\alpha_2=(1,2),\ldots,\alpha_i-\alpha_k=(i,k).
\end{gathered}
\tag{8}
\end{equation}
If, in the form (2), we make the linear substitution
\begin{equation}
 x=\alpha x'+\beta y',
 \qquad
 y=\gamma x'+\delta y'
\tag{9}
\end{equation}
with non-zero determinant
\begin{equation}
 r=\alpha\delta-\beta\gamma,
\tag{10}
\end{equation}
then the form $f(x,y)$ passes into a form $F(x',y')$ of degree $n$, and
\begin{equation}
 F(x',y')=a'_0x'^n+a'_1x'^{n-1}y'+a'_2x'^{n-2}y'^2+\cdots+a'_ny'^n.
\tag{11}
\end{equation}
The coefficients $a'_0,a'_1,\ldots,a'_n$ of the transformed form are obtained as linear and homogeneous expressions in the coefficients $a_0,a_1,\ldots,a_n$, and as homogeneous expressions of degree $n$ in the substitution coefficients $\alpha,\beta,\gamma,\delta$. Thus, by setting $x'=1,y'=0$ or $x'=0,y'=1$, one obtains
\begin{equation}
 a'_0=f(\alpha,\gamma),
 \qquad
 a'_n=f(\beta,\delta).
\tag{12}
\end{equation}
The remaining coefficients $a'$ can be expressed through the derivatives of the function $f(x,y)$, for example
\[
 (n-1)a'_1=\beta f'(\alpha)+\delta f'(\gamma),
\]
which we shall not carry out further.

If we apply the transformation (9) simultaneously to the two pairs of variables $x,y$ and $\alpha_k,\beta_k$, then we obtain the transformed linear factors of $f(x,y)$, namely
\begin{equation}
 (x'\beta'_k)=r(x\beta_k),
\tag{13}
\end{equation}
as follows either from the multiplication rule for determinants or by direct calculation. These factors $(x\beta_k)$ are therefore covariants of $f(x,y)$, though irrational ones, since they do not depend rationally on the coefficients of $f(x)$. Similarly one obtains
\begin{equation}
 (\alpha'_h\beta'_k)=r(\alpha_h\beta_k),
\tag{14}
\end{equation}
according to which the determinants $(\alpha_h\beta_k)$ are to be regarded as irrational invariants. We shall now explain how rational invariants and covariants may be formed from them.

For this purpose it is best to consider the linear transformation in the non-homogeneous form
\begin{equation}
 x=\frac{\alpha x'+\beta}{\gamma x'+\delta}
\tag{15}
\end{equation}
and to apply it to the function $f(x)$ in (1). Then, if
\[
 F(x')=a'_0x'^n+a'_1x'^{n-1}+a'_2x'^{n-2}+\cdots+a'_n
\]
is put, one obtains
\begin{equation}
 F(x')=(\gamma x'+\delta)^n f(x).
\tag{16}
\end{equation}
We also write this equation as
\[
 F(x')=(\alpha x'+\beta)^n
 \left(a_0+a_1\frac1x+a_2\frac1{x^2}+\cdots\right),
\]
from which, if $x'=-\delta:\gamma$, so that $x=\infty$, one obtains
\[
 r^n a_0=(-\gamma)^n F\left(-\frac{\delta}{\gamma}\right).
\]
Or, if $\alpha'_1,\alpha'_2,\ldots,\alpha'_n$ depend on $\alpha_1,\alpha_2,\ldots,\alpha_n$ in the same way as $x'$ depends on $x$, so that
\[
 F(x')=a'_0(x'-\alpha'_1)(x'-\alpha'_2)\cdots(x'-\alpha'_n)
\]
is put, it follows that
\begin{equation}
 r^n a_0=a'_0(\gamma\alpha'_1+\delta)(\gamma\alpha'_2+\delta)\cdots(\gamma\alpha'_n+\delta).
\tag{17}
\end{equation}

Now, by the transformation (15),
\begin{equation}
\begin{aligned}
 x-\alpha_i&=\frac{r(x'-\alpha'_i)}{(\gamma x'+\delta)(\gamma\alpha'_i+\delta)},\\
 \alpha_i-\alpha_k&=\frac{r(\alpha'_i-\alpha'_k)}{(\gamma\alpha'_i+\delta)(\gamma\alpha'_k+\delta)}.
\end{aligned}
\tag{18}
\end{equation}
Now form a product $P$ from any of the factors
\[
 (0,1),(0,2),\ldots,(0,n),
\]
\[
 (1,2),(1,3),\ldots,(n-1,n),
\]
but in such a way that each of the indices $1,2,\ldots,n$ occurs altogether equally often, say $\mu$ times. The index $0$, that is, the variable $x$, is to occur $\nu$ times, and the total number of factors is to be $q$. Since each factor contains two indices, one has
\begin{equation}
 \nu+n\mu=2q.
\tag{19}
\end{equation}
Then from (17) and (18), if $P'$ is obtained from $P$ by replacing $x,\alpha_i$ by $x',\alpha'_i$, one obtains
\begin{equation}
 a_0'{}^\mu P'=a_0^\mu r^{n\mu-q}(\gamma x'+\delta)^\nu P.
\tag{20}
\end{equation}
We now form
\begin{equation}
 C(x,y,a)=a_0^\mu y^\nu\sum P\left(\frac{x}{y}\right),
\tag{21}
\end{equation}
where the sum is taken over all values obtained from $P$ by permuting the indices $1,2,3,\ldots,n$ with one another in all possible ways. This sum is then a symmetric function of the roots of (1), and may therefore be expressed as an integral rational function of the ratios $a_1:a_0,a_2:a_0,\ldots,a_n:a_0$, with no term exceeding the $\mu$th order. Thus $C(x,y,a)$ is an integral homogeneous function of $a_0,a_1,\ldots,a_n$ of order $\mu$, and an integral homogeneous function of $x,y$ of order $\nu$. Since, by (19), (20), it satisfies the condition
\begin{equation}
 C(x',y',a')=r^\lambda C(x,y,a),
\tag{22}
\end{equation}
\begin{equation}
 \lambda=n\mu-q,
 \qquad
 2\lambda=n\mu-\nu,
\tag{23}
\end{equation}
it is a covariant, and in the special case where $\nu=0$ an invariant.

The weight $\lambda$ is, as the second expression shows, always positive or zero, since $n\mu$ is at least equal to $\nu$. The extreme value $n\mu=\nu$, or $\lambda=0$, occurs only when every one of the indices $1,2,3,\ldots,n$ occurs in $P$ only in connection with $0$, that is, when $P$ is a power of $f(x)$ itself.

\end{document}
